\documentclass[a4paper,12pt]{article}

% ---------- Кодировка и язык ----------
\usepackage[T2A]{fontenc}
\usepackage[utf8]{inputenc}
\usepackage[russian]{babel}

% ---------- Математика ----------
\usepackage{amsmath,amssymb,amsthm,mathtools}
\usepackage{icomma}

% ---------- Шрифт и типографика ----------
\usepackage{helvet}
\renewcommand{\familydefault}{\sfdefault}
\usepackage{microtype}
\usepackage{setspace}
\usepackage{parskip}

% ---------- Страница ----------
\usepackage[
  a4paper,
  left=18mm,
  right=18mm,
  top=20mm,
  bottom=20mm
]{geometry}

% ---------- Графика ----------
\usepackage{tikz}
\usetikzlibrary{calc}
\usepackage{tkz-euclide}
\usepackage{pgfplots}
\pgfplotsset{compat=1.18}

% ---------- Таблицы ----------
\usepackage{tabularx,booktabs,longtable,array}

% ---------- Списки ----------
\usepackage{enumitem}

% ---------- Служебное ----------
\usepackage{xcolor}
\usepackage{hyperref}
\usepackage{fancyhdr}

% ---------- Палитра ----------
\definecolor{accent}{HTML}{008080}
\definecolor{darkaccent}{HTML}{004D4D}
\definecolor{textgray}{HTML}{333333}
\definecolor{softgray}{HTML}{6B7373}
\definecolor{lightaccent}{HTML}{E8F2F2}

% ---------- Общие настройки ----------
\onehalfspacing
\setlength{\parindent}{0pt}
\setlength{\headheight}{15pt}
\setlength{\tabcolsep}{6pt}
\renewcommand{\arraystretch}{1.28}
\emergencystretch=2em
\raggedbottom

\hypersetup{
  colorlinks=true,
  linkcolor=darkaccent,
  urlcolor=darkaccent,
  citecolor=darkaccent,
  pdfauthor={Лёвин Артём Александрович},
  pdftitle={Чек-лист. Основы стереометрии: обозначения и первые две аксиомы}
}

\setlist[itemize]{
  leftmargin=7mm,
  itemsep=2pt,
  topsep=3pt
}

\setlist[enumerate,1]{
  leftmargin=8mm,
  itemsep=5pt,
  topsep=4pt,
  label=\arabic*.
}

\setlist[enumerate,2]{
  leftmargin=8mm,
  itemsep=2pt,
  topsep=2pt,
  label=\arabic*)
}

% ---------- Колонтитулы ----------
\pagestyle{fancy}
\fancyhf{}
\fancyhead[L]{\small\textcolor{softgray}{Основы стереометрии}}
\fancyhead[R]{\small\textcolor{softgray}{Ксения Клыкова}}
\fancyfoot[C]{\small\textcolor{softgray}{\thepage}}
\renewcommand{\headrulewidth}{0.3pt}
\renewcommand{\headrule}{
  \hbox to\headwidth{\color{accent!45}\leaders\hrule height \headrulewidth\hfill}
}

% ---------- Вспомогательные команды ----------
\newcommand{\term}[1]{\textcolor{darkaccent}{\textbf{#1}}}
\newcommand{\important}[1]{\textbf{\textcolor{darkaccent}{#1}}}
\newcommand{\plane}[1]{(#1)}
\newcommand{\Axiom}[1]{\textcolor{darkaccent}{\textbf{Аксиома #1.}}}
\newcommand{\solutionmark}{\hfill$\square$}

\begin{document}

% =========================================================
% ТИТУЛЬНЫЙ ЛИСТ
% =========================================================

\begin{titlepage}
\thispagestyle{empty}

\centering

\vspace*{25mm}

{\Large\textcolor{softgray}{ЧЕК-ЛИСТ}\par}

\vspace{10mm}

{\fontsize{26}{31}\selectfont
\bfseries
\textcolor{darkaccent}{
Основы стереометрии
}\par}

\vspace{3mm}

{\Large
Обозначения и первые две аксиомы
\par}

\vspace{14mm}

{\color{accent!70}\rule{0.46\linewidth}{0.8pt}\par}

\vspace{14mm}

{\large
Краткая система обозначений,\\
условия задания плоскости,\\
работа с принадлежностью и доказательствами
\par}

\vfill

{\large
\textbf{Лёвин Артём Александрович}\\[2mm]
эксклюзивно для Ксении Клыковой
\par}

\vspace{8mm}

{\large \today\par}

\vfill

\begin{tikzpicture}[x=1cm,y=1cm]
  \draw[
    accent!55,
    line width=1.1pt
  ]
  (0,0)
  .. controls (2.2,0.55) and (4.1,0.65) .. (6.0,0.18)
  .. controls (8.2,-0.35) and (10.2,-0.15) .. (12.0,0.28)
  .. controls (13.3,0.58) and (14.2,0.38) .. (15.0,0.08);
\end{tikzpicture}

\vspace{4mm}

\end{titlepage}

% =========================================================
% ОСНОВНОЙ ЧЕК-ЛИСТ
% =========================================================

\section*{\textcolor{darkaccent}{1. Язык стереометрии}}

Стереометрия изучает геометрические объекты в пространстве. На первом этапе особенно важно уверенно различать \term{точку}, \term{прямую} и \term{плоскость} и правильно записывать отношения между ними.

\subsection*{\textcolor{accent}{Основные обозначения}}

\begin{table}[h!]
\centering
\caption{Базовые объекты и их обозначения}
\begin{tabularx}{\linewidth}{@{}p{0.20\linewidth}p{0.28\linewidth}X@{}}
\toprule
\textbf{Объект} & \textbf{Обозначение} & \textbf{Комментарий} \\
\midrule
Точка
&
$A$, $B$, $C$, $M$
&
Обозначается заглавной латинской буквой.
\\

Прямая
&
$a$, $b$, $c$ или $AB$
&
Обычно используется строчная латинская буква. Прямую также можно задавать двумя различными точками.
\\

Плоскость
&
$\alpha$, $\beta$, $\gamma$ или $\plane{ABC}$
&
Обычно обозначается строчной греческой буквой. Запись $\plane{ABC}$ допустима, если точки $A$, $B$, $C$ не лежат на одной прямой.
\\
\bottomrule
\end{tabularx}
\end{table}

\subsection*{\textcolor{accent}{Как изображать плоскость}}

Плоскость бесконечна. На чертеже показывается только условный конечный фрагмент плоскости. Чаще всего используется параллелограмм: он помогает передать пространственное расположение объекта.

\begin{center}
\begin{tikzpicture}[scale=0.9]
  \tkzSetUpLine[line width=0.9pt,color=accent]
  \tkzSetUpPoint[color=darkaccent,fill=darkaccent,size=3]
  \tkzSetUpLabel[color=black,font=\small]

  \tkzDefPoint(0,0){P}
  \tkzDefPoint(5.4,0.4){Q}
  \tkzDefPoint(6.5,2.8){R}
  \tkzDefPoint(1.1,2.4){S}

  \tkzDefPoint(2.0,0.85){A}
  \tkzDefPoint(4.2,1.15){B}
  \tkzDefPoint(3.2,2.0){C}
  \tkzDefPoint(6.15,2.25){L}

  \tkzFillPolygon[fill=accent!7](P,Q,R,S)
  \tkzDrawSegments(P,Q Q,R R,S S,P)

  \tkzDrawPoints(A,B,C)
  \tkzLabelPoint[below](A){$A$}
  \tkzLabelPoint[below right](B){$B$}
  \tkzLabelPoint[above](C){$C$}
  \tkzLabelPoint[right](L){$\alpha$}
\end{tikzpicture}
\end{center}

\important{Важно:} форма нарисованного фрагмента плоскости не является её свойством. Параллелограмм, овал или иной аккуратный контур служат только условным изображением бесконечной плоскости.

% ---------------------------------------------------------

\section*{\textcolor{darkaccent}{2. Принадлежность и пересечение}}

\subsection*{\textcolor{accent}{Точка и прямая}}

Если точка $A$ лежит на прямой $a$, записывают
\[
A\in a.
\]

Если точка $C$ на прямой $a$ не лежит:
\[
C\notin a.
\]

\subsection*{\textcolor{accent}{Точка и плоскость}}

Если точка $A$ лежит в плоскости $\alpha$:
\[
A\in\alpha.
\]

Если точка $D$ не лежит в этой плоскости:
\[
D\notin\alpha.
\]

\subsection*{\textcolor{accent}{Прямая и плоскость}}

Если вся прямая $a$ лежит в плоскости $\alpha$, записывают
\[
a\subset\alpha.
\]

Знаки $\in$ и $\subset$ имеют разный смысл:

\begin{itemize}
  \item $A\in\alpha$ --- \textbf{точка} принадлежит плоскости;
  \item $a\subset\alpha$ --- \textbf{прямая целиком} лежит в плоскости.
\end{itemize}

\important{Типичная ошибка:}
\[
A\subset\alpha
\]
для точки $A$ писать не следует. Для точки используется знак $\in$.

\subsection*{\textcolor{accent}{Пересечение}}

Если прямая $a$ имеет с плоскостью $\alpha$ ровно одну общую точку $A$, удобно записать
\[
a\cap\alpha=\{A\}.
\]

Если две прямые пересекаются в точке $J$:
\[
a\cap b=\{J\}.
\]

Если две различные плоскости пересекаются по прямой $AB$:
\[
\alpha\cap\beta=AB.
\]

Если две плоскости совпадают:
\[
\alpha=\beta.
\]

Если они различны:
\[
\alpha\ne\beta.
\]

% =========================================================

\section*{\textcolor{darkaccent}{3. Первая аксиома стереометрии}}

\Axiom{1}
\important{Через любые три точки, не лежащие на одной прямой, проходит плоскость, и притом только одна.}

Иначе говоря, три \term{неколлинеарные} точки однозначно задают плоскость.

\[
A,\ B,\ C\ \text{не лежат на одной прямой}
\quad\Longrightarrow\quad
\exists!\,\plane{ABC}.
\]

Знак $\exists!$ означает: \textit{существует ровно один объект}.

\begin{center}
\begin{tikzpicture}[scale=0.9]
  \tkzSetUpLine[line width=0.9pt,color=accent]
  \tkzSetUpPoint[color=darkaccent,fill=darkaccent,size=3]
  \tkzSetUpLabel[color=black,font=\small]

  \tkzDefPoint(0,0){P}
  \tkzDefPoint(5.3,0.45){Q}
  \tkzDefPoint(6.2,2.8){R}
  \tkzDefPoint(0.9,2.35){S}

  \tkzDefPoint(1.65,0.8){A}
  \tkzDefPoint(4.55,1.0){B}
  \tkzDefPoint(3.0,2.0){C}
  \tkzDefPoint(5.85,2.25){L}

  \tkzFillPolygon[fill=accent!7](P,Q,R,S)
  \tkzDrawSegments(P,Q Q,R R,S S,P)

  \tkzDrawSegments[thin,color=softgray](A,B B,C C,A)

  \tkzDrawPoints(A,B,C)
  \tkzLabelPoint[below left](A){$A$}
  \tkzLabelPoint[below right](B){$B$}
  \tkzLabelPoint[above](C){$C$}
  \tkzLabelPoint[right](L){$\alpha$}
\end{tikzpicture}
\end{center}

\subsection*{\textcolor{accent}{Что означает «не лежат на одной прямой»}}

Точки $A$, $B$, $C$ должны быть расположены так, чтобы \textbf{не существовало одной прямой}, содержащей сразу все три точки.

При этом любые две различные точки всегда определяют некоторую прямую. Это не мешает применению аксиомы.

\subsection*{\textcolor{accent}{Критически важное уточнение}}

Если три точки $A$, $B$, $C$ лежат на одной прямой, неверно говорить, что через них нельзя провести плоскость.

На самом деле:

\[
A,B,C\in a
\quad\Longrightarrow\quad
\text{через }A,B,C\text{ проходит бесконечно много плоскостей}.
\]

Все эти плоскости содержат одну и ту же прямую $a$.

Следовательно, три коллинеарные точки \important{не задают плоскость однозначно}.

% ---------------------------------------------------------

\subsection*{\textcolor{accent}{Пример 1. Вершины треугольника}}

Дан треугольник $ABC$.

По определению его вершины не лежат на одной прямой. Поэтому по первой аксиоме существует единственная плоскость
\[
\plane{ABC}.
\]

Следовательно, вершины любого невырожденного треугольника однозначно задают плоскость.

% ---------------------------------------------------------

\subsection*{\textcolor{accent}{Пример 2. Точка на стороне треугольника}}

Пусть $M\in AB$, $M\ne A$, а $ABC$ --- треугольник.

Точки $A$ и $M$ лежат на прямой $AB$, но точка $C$ на этой прямой не лежит. Поэтому точки $A$, $M$, $C$ неколлинеарны.

Следовательно, существует единственная плоскость
\[
\plane{AMC}.
\]

% ---------------------------------------------------------

\subsection*{\textcolor{accent}{Пример 3. Совпадение двух плоскостей}}

Пусть плоскости $\alpha$ и $\beta$ содержат одни и те же три неколлинеарные точки:
\[
A,B,C\in\alpha,
\qquad
A,B,C\in\beta.
\]

По первой аксиоме через $A$, $B$, $C$ проходит \textbf{единственная} плоскость.

Следовательно,
\[
\boxed{\alpha=\beta}.
\]

\important{Ключ к рассуждению:} слово «единственная» в первой аксиоме позволяет доказывать совпадение плоскостей.

% =========================================================

\section*{\textcolor{darkaccent}{4. Доказательство от противного}}

В задачах по стереометрии часто удобно использовать \term{метод доказательства от противного}.

\subsection*{\textcolor{accent}{Алгоритм}}

\begin{enumerate}
  \item Предположите противоположное тому, что требуется доказать.
  \item Выведите из предположения необходимые следствия.
  \item Получите противоречие с условием задачи, определением или аксиомой.
  \item Сделайте вывод, что первоначальное предположение было неверным.
\end{enumerate}

\subsection*{\textcolor{accent}{Пример}}

Даны точки $A$, $B$, $C$, $D$, причём никакие три из них не лежат на одной прямой и
\[
D\notin\plane{ABC}.
\]

Докажем:
\[
\plane{ABC}\ne\plane{ABD}.
\]

\textbf{Доказательство.}

Предположим противное:
\[
\plane{ABC}=\plane{ABD}.
\]

Но точка $D$ принадлежит плоскости $\plane{ABD}$. При совпадении плоскостей тогда должно выполняться
\[
D\in\plane{ABC}.
\]

Это противоречит условию
\[
D\notin\plane{ABC}.
\]

Следовательно,
\[
\boxed{\plane{ABC}\ne\plane{ABD}}.
\]
\solutionmark

% =========================================================

\section*{\textcolor{darkaccent}{5. Вторая аксиома стереометрии}}

\Axiom{2}
\important{Если две точки прямой принадлежат плоскости, то вся прямая принадлежит этой плоскости.}

Пусть $A$ и $B$ --- две различные точки прямой $a$. Если
\[
A\in\alpha,
\qquad
B\in\alpha,
\]
то
\[
\boxed{a\subset\alpha}.
\]

Если прямая задана точками $A$ и $B$, можно записать:
\[
A\in\alpha,\quad B\in\alpha
\quad\Longrightarrow\quad
AB\subset\alpha.
\]

\begin{center}
\begin{tikzpicture}[scale=0.9]
  \tkzSetUpLine[line width=0.9pt,color=accent]
  \tkzSetUpPoint[color=darkaccent,fill=darkaccent,size=3]
  \tkzSetUpLabel[color=black,font=\small]

  \tkzDefPoint(0,0){P}
  \tkzDefPoint(5.5,0.4){Q}
  \tkzDefPoint(6.4,2.9){R}
  \tkzDefPoint(0.9,2.5){S}

  \tkzDefPoint(0.8,0.72){U}
  \tkzDefPoint(2.0,1.20){A}
  \tkzDefPoint(4.0,2.00){B}
  \tkzDefPoint(5.6,2.64){V}
  \tkzDefPoint(6.02,2.27){L}

  \tkzFillPolygon[fill=accent!7](P,Q,R,S)
  \tkzDrawSegments(P,Q Q,R R,S S,P)

  \tkzDrawSegment[line width=1.2pt,color=darkaccent](U,V)

  \tkzDrawPoints(A,B)
  \tkzLabelPoint[below](A){$A$}
  \tkzLabelPoint[above](B){$B$}
  \tkzLabelPoint[right](L){$\alpha$}
\end{tikzpicture}
\end{center}

\subsection*{\textcolor{accent}{Следствие для любой точки прямой}}

Если уже доказано, что
\[
a\subset\alpha,
\]
то для любой точки $C\in a$ автоматически выполняется
\[
C\in\alpha.
\]

Схема рассуждения:
\[
A,B\in\alpha
\overset{\text{аксиома 2}}{\Longrightarrow}
a\subset\alpha,
\qquad
C\in a
\Longrightarrow
C\in\alpha.
\]

\subsection*{\textcolor{accent}{Пример}}

Пусть
\[
A\in\alpha,\qquad B\in\alpha,\qquad C\in AB.
\]

Тогда по аксиоме 2:
\[
AB\subset\alpha.
\]

Поскольку $C\in AB$, получаем
\[
\boxed{C\in\alpha}.
\]

% =========================================================

\section*{\textcolor{darkaccent}{6. Как связывать две аксиомы}}

В задачах часто требуется использовать аксиомы последовательно.

\subsection*{\textcolor{accent}{Типовая схема}}

Пусть $ABC$ --- треугольник, $M\in AB$.

\begin{enumerate}
  \item Точки $A$, $B$, $C$ неколлинеарны, поэтому по аксиоме 1 существует единственная плоскость $\plane{ABC}$.
  \item Точки $A$ и $B$ принадлежат этой плоскости:
  \[
  A,B\in\plane{ABC}.
  \]
  \item По аксиоме 2:
  \[
  AB\subset\plane{ABC}.
  \]
  \item Так как $M\in AB$, получаем:
  \[
  M\in\plane{ABC}.
  \]
  \item Точки $A$, $M$, $C$ неколлинеарны, поэтому они определяют единственную плоскость.
  \item Плоскость $\plane{ABC}$ уже содержит все три точки $A$, $M$, $C$. Следовательно,
  \[
  \boxed{\plane{AMC}=\plane{ABC}}.
  \]
\end{enumerate}

% =========================================================

\section*{\textcolor{darkaccent}{7. Типичные ошибки}}

\begin{enumerate}
  \item \textbf{Перепутать $\in$ и $\subset$.}

  Для точки:
  \[
  A\in\alpha.
  \]

  Для прямой:
  \[
  a\subset\alpha.
  \]

  \item \textbf{Забыть условие неколлинеарности в аксиоме 1.}

  Трёх точек недостаточно. Требуется, чтобы они не лежали на одной прямой.

  \item \textbf{Считать, что три коллинеарные точки не принадлежат никакой плоскости.}

  Они принадлежат бесконечному числу плоскостей, содержащих их общую прямую.

  \item \textbf{Считать, что одной точки прямой достаточно для принадлежности всей прямой плоскости.}

  Из
  \[
  A\in a,\qquad A\in\alpha
  \]
  нельзя сделать вывод
  \[
  a\subset\alpha.
  \]

  Для применения аксиомы 2 нужны \important{две различные точки} прямой.

  \item \textbf{Не ссылаться на аксиому в доказательстве.}

  Лучше писать:
  \[
  A,B\in\alpha
  \overset{\text{аксиома 2}}{\Longrightarrow}
  AB\subset\alpha.
  \]

  \item \textbf{Путать рисунок с математическим фактом.}

  По расположению объектов на плоском чертеже нельзя делать вывод о принадлежности или пересечении без условия задачи либо доказательства.

  \item \textbf{Использовать доказательство от противного без явного предположения.}

  Начинайте формулировкой:
  \[
  \text{«Предположим противное: \ldots»}
  \]
\end{enumerate}

% =========================================================

\section*{\textcolor{darkaccent}{8. Быстрый алгоритм перед решением задачи}}

Перед доказательством задайте себе пять вопросов.

\begin{enumerate}
  \item Какие объекты даны: точки, прямые, плоскости?
  \item Какие принадлежности известны?
  \item Есть ли три неколлинеарные точки, позволяющие применить аксиому 1?
  \item Есть ли две точки одной прямой, принадлежащие плоскости, чтобы применить аксиому 2?
  \item Нужно ли доказать равенство или различие плоскостей? Поможет ли единственность плоскости или метод от противного?
\end{enumerate}

\subsection*{\textcolor{accent}{Минимум, который нужно помнить наизусть}}

\begin{table}[h!]
\centering
\caption{Главные факты занятия}
\begin{tabularx}{\linewidth}{@{}p{0.23\linewidth}X@{}}
\toprule
\textbf{Элемент} & \textbf{Формулировка} \\
\midrule
Аксиома 1
&
Через три точки, не лежащие на одной прямой, проходит единственная плоскость.
\\

Аксиома 2
&
Если две точки прямой принадлежат плоскости, то вся прямая принадлежит этой плоскости.
\\

Точка в плоскости
&
$A\in\alpha$.
\\

Прямая в плоскости
&
$a\subset\alpha$.
\\

Совпадение плоскостей
&
Если две плоскости содержат одни и те же три неколлинеарные точки, они совпадают.
\\

Метод от противного
&
Предположить отрицание требуемого утверждения и получить противоречие.
\\
\bottomrule
\end{tabularx}
\end{table}

% =========================================================
% ТРЕНИРОВОЧНЫЙ БЛОК
% =========================================================

\newpage

\section*{\textcolor{darkaccent}{Тренировочный блок --- Путь А}}

\textcolor{softgray}{
10 упражнений по нарастающей сложности. Решения сначала выполните самостоятельно.
}

\subsection*{\textcolor{accent}{Средний уровень}}

\begin{enumerate}
  \item
  Точки $A$, $B$, $C$ не лежат на одной прямой.

  Сколько плоскостей проходит через все три точки? Назовите эту плоскость.

  \item
  Дан треугольник $PQR$.

  Можно ли однозначно задать плоскость его вершинами? Обоснуйте ответ ссылкой на аксиому.

  \item
  Точки $A$ и $B$ принадлежат плоскости $\alpha$, причём $A\ne B$.

  Какой вывод можно сделать о прямой $AB$?
\end{enumerate}

\subsection*{\textcolor{accent}{Уровень выше среднего}}

\begin{enumerate}[resume]
  \item
  Точки $A$, $B$, $C$ лежат на одной прямой.

  Сколько плоскостей может проходить одновременно через все три точки?

  \item
  Известно:
  \[
  A\in\alpha,\qquad
  B\in\alpha,\qquad
  M\in AB.
  \]

  Докажите:
  \[
  M\in\alpha.
  \]

  \item
  Плоскости $\alpha$ и $\beta$ содержат точки $P$, $Q$, $R$, которые не лежат на одной прямой.

  Докажите:
  \[
  \alpha=\beta.
  \]

  \item
  Дан треугольник $ABC$, точка $M$ лежит на стороне $AB$, причём $M\ne A$.

  Докажите:
  \[
  \plane{AMC}=\plane{ABC}.
  \]

  \item
  Даны четыре точки $A$, $B$, $C$, $D$. Никакие три из них не лежат на одной прямой и
  \[
  D\notin\plane{ABC}.
  \]

  Докажите:
  \[
  \plane{ABC}\ne\plane{ABD}.
  \]

  \item
  Плоскости $\alpha$ и $\beta$ различны. Известно, что
  \[
  A\in\alpha,\qquad
  B\in\alpha,\qquad
  A\in\beta,\qquad
  B\in\beta,
  \qquad A\ne B.
  \]

  Что можно утверждать о прямой $AB$?
\end{enumerate}

\subsection*{\textcolor{accent}{Сложное упражнение}}

\begin{enumerate}[resume]
  \item
  Точки $A$, $B$, $C$ не лежат на одной прямой, а
  \[
  D\notin\plane{ABC}.
  \]

  Точка $M$ лежит на прямой $CD$, причём
  \[
  M\ne C,\qquad M\ne D.
  \]

  Докажите методом от противного:
  \[
  M\notin\plane{ABC}.
  \]
\end{enumerate}
\end{document}